\documentclass[11pt,letterpaper]{article}
\usepackage[T1]{fontenc}
\usepackage[utf8]{inputenc}
\usepackage{lmodern, microtype, amsmath,amssymb,amsthm,mathtools,booktabs,enumitem}
\usepackage[margin=0.91in]{geometry}
\usepackage[colorlinks=true,urlcolor=blue,linkcolor=black,citecolor=black]{hyperref}
\usepackage{fancyhdr}
\setlength{\headheight}{14pt}
\pagestyle{fancy}\fancyhf{}\lhead{From Prephysical Hypotheses to Physical Models}\rhead{Mathematical supplement}\cfoot{\thepage}
\newtheorem{theorem}{Theorem}
\newtheorem{proposition}[theorem]{Proposition}
\newtheorem{corollary}[theorem]{Corollary}
\newtheorem{remark}[theorem]{Remark}
\newcommand{\D}{\mathcal{D}}
\newcommand{\Ssep}{\mathcal{S}}
\newcommand{\K}{\mathcal{K}}
\newcommand{\tr}{\operatorname{Tr}}
\newcommand{\id}{\mathbf{1}}
\newcommand{\C}{\mathbb{C}}
\title{\textbf{Composite-State Selection and Multipartite Closure:}\\
An Operational Countermodel and a Conditional Entangling-Gate Criterion}
\author{Kirtus Justus\\\small Research direction and scholarly responsibility\\
\small Port Coquitlam, British Columbia, Canada\\[4pt]
\small AI-assisted internal mathematical checking: ChatGPT (GPT-6), Director G}
\date{12 November 2026, B.C. Pacific Time\\[2pt]
\small Initial External Grounding; research stage OPEN; no ontology selected}
\begin{document}
\maketitle
\begin{abstract}
We distinguish three logical burdens often collapsed in attempts to infer a complete
quantum theory from single-system pure-state geometry: selecting a two-system state
space, preserving that selection under independent multipartite preparation, and
justifying entangling dynamics. A continuous family of stipulated two-system
operational models has exactly the standard elementary complex two-level state
spaces and measurements, and is locally tomographic, but differs in its allowed
joint states. A proposed uniform four-system extension fails an exact parallel-
preparation test by $t(1-t)/4$ for every $0<t<1$. Conversely, assuming the standard
complex operator representation, pure product preparations, all local unitaries,
convex mixing and a controlled-not operation is sufficient to force the full
ordinary two-system density-matrix state space. These are conditional mathematical
results, not an origin of complex structure, quantum probabilities, an entangling
gate, dynamics, classical mechanics or physical spacetime. The source research
baseline remains unchanged.
\end{abstract}

\section{Question, epistemic scope and exact definitions}
The original scientific target is an explained emergence of both quantum and classical
mechanics, with a calibrated correspondence from prephysical primitives. The present
note answers a narrower \emph{mathematical falsification question}: does exhaustive
knowledge of an elementary system's conventional complex projective geometry,
local quantum measurements and local reversible symmetries determine composition?
A numerical or analytic answer here is not a selection of a real ontology.

The term ``two-level system'' here refers \emph{only} to a stipulated mathematical
complex Hilbert space $H=\C^2$ with the usual inner product. It is neither an
operationally identified physical qubit nor a qubit derived from prephysical
hypotheses. Its Born probabilities and tensor products are part of this test's
\emph{assumptions}, not conclusions.

Let $\D_n$ be the compact convex set of positive trace-one operators on
$H^{\otimes n}$, and let
\[
\Ssep_n = \operatorname{conv}\{\rho_1\otimes\cdots\otimes\rho_n:
\rho_j\in\D_1\}
\]
be the fully separable subset. For $0\leq t\leq 1$ set
\begin{equation}\label{eq:Kt}
\K_t^{(2)}=(1-t)\Ssep_2+t\D_2
=\{(1-t)\sigma+t\tau:\sigma\in\Ssep_2,\;\tau\in\D_2\}.
\end{equation}
Minkowski addition is intended, not an unmotivated probabilistic estimate of what
nature prefers. At the endpoints $\K_0^{(2)}=\Ssep_2$ and $\K_1^{(2)}=\D_2$.

All tests use standard single-system density matrices, local quantum POVMs and
local completely positive trace-preserving (CPTP) maps, as \emph{given inputs}.
For joint measurements we allow the effects $0\preceq E\preceq\id_4$ and their
probabilities $\tr(E\rho)$. Joint binary effects can be implemented by
measure-and-prepare instruments whose conditional output is a fixed separable
state, so no entangled postmeasurement state is implicitly required. The
\emph{no-restriction hypothesis} for joint effects and arbitrary entangling
unitaries are not assumed.

\section{What is and is not selected at two systems}
\begin{proposition}[Operational pair family]\label{prop:pair}
For every $t\in[0,1]$, $\K_t^{(2)}$ is a compact convex bipartite state space
containing all independent product preparations, stable under arbitrary local
CPTP maps, with well-defined normalized positive quantum measurement effects.
All members have the same elementary preparations and measurement probabilities.
They are locally tomographic: product effects suffice to distinguish their
bipartite states.
\end{proposition}
\begin{proof}
Convexity and compactness follow from those of $\Ssep_2$ and $\D_2$.
Because $\Ssep_2\subseteq\D_2$, every $\sigma\in\Ssep_2$ is in
$\K_t^{(2)}$, by taking $(1-t)\sigma+t\sigma$. The tensor product of two
local CPTP maps sends $\Ssep_2$ to itself and $\D_2$ to itself, hence preserves
$\K_t^{(2)}$. Ordinary positive effects assign numbers in $[0,1]$ to all
allowed states. To distinguish states locally, choose on each factor the four
rank-one effects associated with $|0\rangle,|1\rangle,|+\rangle,
(|0\rangle+i|1\rangle)/\sqrt{2}$. They span the real space of $2\times2$
Hermitian matrices; their sixteen products span the real space of $4\times4$
Hermitian matrices. Thus their expectation values determine the bipartite
operator, regardless of which $\K_t^{(2)}$ contains it. Finally, partial traces
have valid elementary marginals. The two-system model makes no assertion about
unrestricted transformations or a consistent arbitrary-size composite.
\end{proof}

Write $|\Phi\rangle=(|00\rangle+|11\rangle)/\sqrt{2}$ and
$P=|\Phi\rangle\langle\Phi|$. Define the distinguished allowed state
\begin{equation}\label{eq:rho}
\rho_t=(1-t)|00\rangle\langle00|+tP \in\K_t^{(2)}.
\end{equation}

\begin{proposition}[Exact separating effect]\label{prop:bell}
For the physically unestablished but mathematically admissible joint effect $P$,
\[
\sup_{\rho\in\K_t^{(2)}}\tr(P\rho)=\frac{1+t}{2},
\qquad 0\leq t\leq1.
\]
The extremizer $\rho_t$ has partial-transpose minimum eigenvalue $-t/2$.
Under the fixed conventional CHSH settings $A_0=Z$, $A_1=X$,
$B_{0,1}=(Z\pm X)/\sqrt{2}$, it yields
$\langle\mathrm{CHSH}\rangle_{\rho_t}=\sqrt{2}(1+t)$, which exceeds $2$
when $t>\sqrt{2}-1$.
\end{proposition}
\begin{proof}
For any normalized product vector $|a\rangle|b\rangle$,
$|\langle\Phi|a\otimes b\rangle|^2\leq1/2$ by Cauchy--Schwarz, and
convexity extends the inequality to $\Ssep_2$. The maximum of $\tr(P\tau)$
over $\D_2$ is one. Taking the weighted sum gives the upper bound
$(1-t)/2+t$; $|00\rangle\langle00|$ and $P$ attain the constituent bounds.
The partial transpose of \eqref{eq:rho} has spectrum
$\{1-t/2,t/2,t/2,-t/2\}$ (with zeros appropriately combined at $t=0$).
The CHSH expression follows from $\mathrm{CHSH}=\sqrt{2}(Z\otimes Z+X\otimes X)$
for the specified settings. \end{proof}

In particular, even an observed CHSH violation at these settings would not,
without further premises, select the entire two-system density-matrix set.
This observation remains within the constructed family; no empirical CHSH
measurement was performed here.

\section{An exact obstruction at four systems}
One might attempt to use the \emph{same} interpolation parameter at every
system size:
\begin{equation}\label{eq:naive}
\K_t^{(4)}\stackrel{?}{=}(1-t)\Ssep_4+t\D_4.
\end{equation}
Independent preparation of any two allowed pair states would require
$\rho\otimes\sigma\in\K_t^{(4)}$ for all
$\rho,\sigma\in\K_t^{(2)}$. This closure condition is separate from local
tomography and from closure under local channels.

\begin{theorem}[Uniform extension fails parallel-preparation closure]
\label{thm:four}
For every $0<t<1$, the proposed four-party state set \eqref{eq:naive}
excludes an independent product of two states permitted by \eqref{eq:Kt}.
The violation of an exactly sharp positive-effect bound is
\begin{equation}\label{eq:gap}
\Delta(t)=\frac{t(1-t)}4>0.
\end{equation}
At $t=1/2$ this is exactly $1/16$.
\end{theorem}
\begin{proof}
Let $P_{12}=P\otimes\id_4$ and $P_{34}=\id_4\otimes P$,
acting on disjoint elementary pairs in $H^{\otimes4}$. They commute. Define
\[
E=\id_{16}-(\id_{16}-P_{12})(\id_{16}-P_{34})
=\id_{16}-(\id_4-P)\otimes(\id_4-P).
\]
This is a projector, so it is a legitimate effect ($0\preceq E\preceq\id$).
On each fully product pure state the two pair fidelities $f_1,f_2$ are at most
$1/2$. Hence $\tr(E\omega)=1-(1-f_1)(1-f_2)\leq3/4$.
Convexity gives this bound for all $\Ssep_4$; it is attained by $|0000\rangle$.
On $\D_4$ the supremum is one, attained by $P\otimes|00\rangle\langle00|$.
Thus the exact support bound over the proposed set is
\begin{equation}\label{eq:support}
\sup_{\omega\in\K_t^{(4)}}\tr(E\omega)=\frac34+\frac t4.
\end{equation}
Independently prepare $\rho_t$ from \eqref{eq:rho} on each of the pairs
$(1,2)$ and $(3,4)$. The pair Bell fidelity is $(1+t)/2$, so the product
preparation has
\[
\tr[E(\rho_t\otimes\rho_t)]
=1-\left(1-\frac{1+t}{2}\right)^2
=1-\frac{(1-t)^2}{4}.
\]
Subtracting \eqref{eq:support} yields $\Delta(t)=t(1-t)/4$.
It is strictly positive for every interior $t$, establishing the failure.
The endpoints pass this particular test.\end{proof}

\begin{remark}[Precisely what the counterexample defeats]
The two-system $\K_t^{(2)}$ remains a valid \emph{pairwise} operational toy
state set. The failed assertion is that the same simple formula defines a
consistent composition rule simultaneously for pairs and four elementary
systems. This theorem does not exclude all alternative multipartite extensions,
nor claim that every intermediate two-system composite is impossible. At
$t=0$, the fully separable hierarchy is closed under arbitrary tensor
products and discarding; at $t=1$ the usual quantum hierarchy is similarly
closed. Those two controls are retained.
\end{remark}

\section{A conditional selection principle, and its cost}
\begin{theorem}[Specified entangling gate forces the usual two-system cone]
\label{thm:cnot}
Let $\K\subseteq\D_2$ be a convex set containing all pure product
preparations. Assume (a) invariance under arbitrary local unitaries
$U\otimes V$ and (b) closure under the exact controlled-not operation
$\operatorname{CNOT}|a,b\rangle=|a,a\oplus b\rangle$.
Then $\K=\D_2$.
\end{theorem}
\begin{proof}
Any normalized $|\psi\rangle\in H\otimes H$ has a Schmidt decomposition
\[
|\psi\rangle=(U\otimes V)
\bigl(s_0|00\rangle+s_1|11\rangle\bigr),
\quad s_0,s_1\geq0,\quad s_0^2+s_1^2=1.
\]
The bracketed state is
$\operatorname{CNOT}[(s_0|0\rangle+s_1|1\rangle)\otimes|0\rangle]$,
where the input is a pure product. By the two closure assumptions, every
pure two-system state belongs to $\K$. Convex mixtures of pure state
projectors exhaust $\D_2$. Since $\K\subseteq\D_2$, equality follows.
\end{proof}

\noindent\textbf{Scope warning.} The theorem takes as inputs the Hilbert-space
embedding, the ordinary complex operator representation of the joint system,
a named entangling gate, and standard notions of purity and convex mixing.
It proves no \emph{prephysical} origin of those objects. A weaker or different
operational entangling principle requires its own proof. Allowing only local
unitaries does \emph{not} suffice. The theorem is a diagnostic for identifying
exactly which additional compositional postulate performs the selection.

\section{Scientific disposition and reproducibility}
\begin{center}
\begin{tabular}{p{0.52\linewidth}p{0.34\linewidth}}
\toprule
Target of assessment & Scope of finding\\\midrule
Nonselection of two-system composition by local geometry and local quantum effects & Exact mathematical countermodels, supported in stated scope\\
Uniform multipartite interpolation \eqref{eq:naive} for $0<t<1$ & Refuted by the positive-effect witness\\
Composition selected when CNOT is assumed on $\K\subseteq\D_2$ & Conditional theorem, not emergent physics\\
Origin of Born probabilities, entangling interactions, clocks and classical mechanics & Not identified by these calculations\\
Total prephysical-to-physical ontology & Not identified; Initial External Grounding stays open\\\bottomrule
\end{tabular}
\end{center}

The independent verification script \texttt{\detokenize{verification/verify_composite_selection.py}}
retains 101 equally spaced $t$ values including both endpoints, 200 seeded
complex-Gaussian target vectors checked through numerical Schmidt synthesis,
and an exact symbolic factorization in SymPy. It confirms the sixteen-dimensional
real span of product preparations/effects, the partial-transpose and CHSH values,
$\Delta(1/2)=1/16$, and the full parameter-grid witness. These are
\emph{deterministic numerical checks and exact algebra}; the grid does not
measure nature, estimate the prevalence of model failure, or supply a
statistical confidence level. The mathematical proofs above, rather than
floating-point tolerances, control the stated universally quantified findings.
No trial is discarded because of an inconvenient value.

The present calculation was produced and self-checked by one AI participant
under the human research direction of Kirtus Justus. It is internal research
provenance, not an independent sub-agent, independent expert review, empirical
validation, or scientific endorsement. No source baseline has been edited.

\begin{thebibliography}{9}\small
\bibitem{Muller2021} M. P. M\"uller, ``Probabilistic theories and reconstructions of quantum theory,''
\emph{SciPost Physics Lecture Notes} \textbf{28} (2021).
\url{https://doi.org/10.21468/SciPostPhysLectNotes.28}.
\bibitem{Barnum2014} H. Barnum, M. P. M\"uller, and C. Ududec,
``Higher-order interference and single-system postulates characterizing quantum theory,''
\emph{New Journal of Physics} \textbf{16}, 123029 (2014).
\url{https://doi.org/10.1088/1367-2630/16/12/123029}.
\end{thebibliography}
\end{document}
